% !TEX program = xelatex
% Lernplatt - by Can Siebert
% Kurs 22 (Bo 143) - Tag 6 - Loesungsheft
% Omnichannel-Strategie & Operationalisierung digitaler Vertriebsstrategien

\documentclass[11pt,a4paper,oneside]{article}

\usepackage{polyglossia}
\setdefaultlanguage{german}

\usepackage[a4paper,margin=2.5cm]{geometry}

\usepackage{fontspec}
\defaultfontfeatures{Ligatures=TeX}

\IfFontExistsTF{Charter}{\setmainfont{Charter}}{%
  \setmainfont{Noto Serif}}
\IfFontExistsTF{Source Sans 3}{\setsansfont{Source Sans 3}}{%
  \IfFontExistsTF{Source Sans Pro}{\setsansfont{Source Sans Pro}}{%
    \setsansfont{Liberation Sans}}}

\newcommand{\caveatfontfallback}{\itshape}
\IfFontExistsTF{Caveat}{\newfontfamily\caveatfont{Caveat}[Scale=1.15]}{\newcommand\caveatfont{\caveatfontfallback}}

\setmonofont{Liberation Mono}

\usepackage{microtype}
\usepackage{eurosym}
\linespread{1.15}

\usepackage{xcolor}
\definecolor{lpInk}{RGB}{20,28,38}
\definecolor{lpAccent}{RGB}{22,90,140}
\definecolor{lpAccentLight}{RGB}{210,228,240}
\definecolor{lpAmber}{RGB}{222,138,30}
\definecolor{lpAmberLight}{RGB}{255,243,217}
\definecolor{lpAmberBorder}{RGB}{196,118,18}
\definecolor{lpGray}{RGB}{96,104,114}
\definecolor{lpGrayLight}{RGB}{238,240,243}
\definecolor{lpGreen}{RGB}{34,120,72}
\definecolor{lpGreenLight}{RGB}{222,238,228}
\definecolor{lpGreenBorder}{RGB}{24,96,58}
\definecolor{lpL1}{RGB}{34,120,72}
\definecolor{lpL2}{RGB}{22,90,140}
\definecolor{lpL3}{RGB}{170,42,52}

\usepackage[most]{tcolorbox}
\tcbuselibrary{breakable,skins}

\newtcolorbox{infobox}[1][Hinweis]{enhanced, breakable, colback=lpAccentLight, colframe=lpAccent, boxrule=0.6pt, arc=2pt, left=10pt,right=10pt,top=8pt,bottom=8pt, fonttitle=\sffamily\bfseries\color{white}, coltitle=white, title={#1}, attach boxed title to top left={xshift=8pt,yshift=-8pt}, boxed title style={size=small, colback=lpAccent, colframe=lpAccent, boxrule=0.4pt, arc=1pt}, top=14pt}

\newtcolorbox{loesungbox}[1][Musterlösung]{enhanced, breakable, colback=lpGreenLight, colframe=lpGreenBorder, boxrule=0.6pt, arc=2pt, left=10pt,right=10pt,top=8pt,bottom=8pt, fonttitle=\sffamily\bfseries\color{white}, coltitle=white, title={#1}, attach boxed title to top left={xshift=8pt,yshift=-8pt}, boxed title style={size=small, colback=lpGreenBorder, colframe=lpGreenBorder, boxrule=0.4pt, arc=1pt}, top=14pt}

\newtcolorbox{merkbox}[1][Managementhinweis]{enhanced, breakable, colback=lpAmberLight, colframe=lpAmberBorder, boxrule=0.6pt, arc=2pt, left=10pt,right=10pt,top=8pt,bottom=8pt, fonttitle=\sffamily\bfseries\color{white}, coltitle=white, title={#1}, attach boxed title to top left={xshift=8pt,yshift=-8pt}, boxed title style={size=small, colback=lpAmberBorder, colframe=lpAmberBorder, boxrule=0.4pt, arc=1pt}, top=14pt}

\usepackage{enumitem}
\setlist{itemsep=2pt, topsep=4pt, parsep=0pt}
\setlist[itemize,1]{label={\textbullet},leftmargin=1.6em}
\setlist[itemize,2]{label={\textendash},leftmargin=1.4em}
\setlist[enumerate,1]{label=\arabic*., leftmargin=1.8em}

\usepackage{tabularx}
\usepackage{array}
\usepackage{booktabs}
\usepackage{longtable}

\usepackage{fancyhdr}
\usepackage{lastpage}
\pagestyle{fancy}
\fancyhf{}
\renewcommand{\headrulewidth}{0.3pt}
\renewcommand{\footrulewidth}{0pt}
\fancyhead[L]{\footnotesize\sffamily\color{lpGray}Lösungsheft \textperiodcentered{} Kurs 22 \textperiodcentered{} Tag 6}
\fancyhead[R]{\footnotesize\sffamily\color{lpGray}Omnichannel \& Operationalisierung}
\fancyfoot[L]{\footnotesize\sffamily\color{lpGray}\textcopyright{} Can Siebert}
\fancyfoot[R]{\footnotesize\sffamily\color{lpGray}\thepage{} / \pageref{LastPage}}

\fancypagestyle{plain}{\fancyhf{}\renewcommand{\headrulewidth}{0pt}\fancyfoot[C]{\footnotesize\sffamily\color{lpGray}\thepage{} / \pageref{LastPage}}}

\usepackage[explicit]{titlesec}
\titleformat{\section}[block]{\sffamily\Large\bfseries\color{lpAccent}}{\thesection.}{0.6em}{#1}[{\vspace{2pt}\color{lpAccent}\titlerule[0.6pt]}]
\titleformat{\subsection}[block]{\sffamily\large\bfseries\color{lpInk}}{\thesubsection}{0.6em}{#1}
\titlespacing*{\section}{0pt}{14pt}{8pt}
\titlespacing*{\subsection}{0pt}{10pt}{4pt}

\usepackage[hidelinks,unicode]{hyperref}
\usepackage{tikz}
\usetikzlibrary{shapes.geometric,arrows.meta,positioning,calc,fit,backgrounds}

\newcommand{\akzent}[1]{{\caveatfont\color{lpAmber}#1}}
\newcommand{\fachbegriff}[1]{\textbf{#1}}

\newcommand{\niveaubadge}[2]{\tcbox[on line, colback=#1, colframe=#1, coltext=white, boxrule=0pt, arc=2pt, left=5pt,right=5pt,top=1pt,bottom=1pt, nobeforeafter]{\sffamily\bfseries\footnotesize #2}}
\newcommand{\nivLi}{\niveaubadge{lpL1}{L1\,\textbullet\,Wissen}}
\newcommand{\nivLii}{\niveaubadge{lpL2}{L2\,\textbullet\,Anwendung}}
\newcommand{\nivLiii}{\niveaubadge{lpL3}{L3\,\textbullet\,Management}}

\newcommand{\zeit}[1]{\tcbox[on line, colback=lpGrayLight, colframe=lpGrayLight, coltext=lpInk, boxrule=0pt, arc=2pt, left=5pt,right=5pt,top=1pt,bottom=1pt, nobeforeafter]{\sffamily\footnotesize\textbf{$\sim$\,#1}}}

\newcommand{\aufgabenkopf}[4]{\par\vspace{6pt}\noindent{\sffamily\Large\bfseries\color{lpAccent}Aufgabe #1 \textendash{} #2}\par\vspace{2pt}\noindent{\color{lpAccent}\rule{\linewidth}{0.6pt}}\par\vspace{4pt}\noindent #3 \hfill \zeit{#4}\par\vspace{6pt}}

\newcommand{\ytquery}[1]{%
  \par\vspace{4pt}
  \noindent{\footnotesize\sffamily\color{lpGray}\textbf{Lernvideo zum Thema:}\ %
    \href{https://studyflix.de/search?query=#1}{\textcolor{lpAccent}{Studyflix-Treffer}}\ ·\ %
    \href{https://www.youtube.com/results?search\_query=#1}{\textcolor{lpAccent}{YouTube-Treffer}}\\%
    \textit{\textcolor{lpGray}{Stichworte: #1}}}\par
}
\newcommand{\ytvideo}[2]{%
  \par\vspace{4pt}
  \noindent{\footnotesize\sffamily\color{lpGray}\textbf{Lernvideo:}\ \href{#2}{\textcolor{lpAccent}{#1}}}\par
}

\begin{document}

\begin{titlepage}
  \pagestyle{empty}
  \vspace*{1.5cm}
  \noindent{\sffamily\footnotesize\color{lpGray}LERNPLATT \textperiodcentered{} BY CAN SIEBERT}\\[2pt]
  \noindent{\color{lpAccent}\rule{\linewidth}{1.2pt}}\\[4pt]
  \noindent{\sffamily\footnotesize\color{lpGray}LÖSUNGSHEFT \textperiodcentered{} MODUL 3 \textperiodcentered{} TAG 6 \textperiodcentered{} KURS 22 (Bo 143) \textperiodcentered{} 5.\,Juni 2026}

  \vspace{3cm}
  {\sffamily\fontsize{30}{34}\selectfont\bfseries\color{lpInk}Lösungsheft\\Tag 6\par}

  \vspace{0.7cm}
  {\sffamily\large\color{lpGray}Musterlösungen zu den elf Aufgaben\\des Aufgabenhefts \textendash{} Omnichannel-\\Strategie \& Operationalisierung.\par}

  \vspace{1.5cm}
  {\caveatfont\LARGE\color{lpGreen}Musterlösungen \textperiodcentered{} nicht die einzige Antwort}\par

  \vfill

  \noindent\begin{minipage}{0.55\linewidth}
    {\sffamily\footnotesize\color{lpGray}Hinweis:}\\
    {\sffamily\small\color{lpInk}Musterlösungen sind Diskussionsbasis,\\nicht die einzige korrekte Lösung.}\\[10pt]
    {\sffamily\footnotesize\color{lpGray}Begleitend:}\\
    {\sffamily\small\color{lpInk}Aufgabenheft \& Lehrskript Tag 6}
  \end{minipage}\hfill
  \begin{minipage}{0.4\linewidth}
    \raggedleft
    {\sffamily\footnotesize\color{lpGray}Dozent:}\\
    {\sffamily\small\color{lpInk}Can Siebert}\\[10pt]
    {\sffamily\footnotesize\color{lpGray}Niveau-Mix:}\\
    {\sffamily\small\color{lpInk}3\,L1 \textperiodcentered{} 5\,L2 \textperiodcentered{} 3\,L3}
  \end{minipage}

  \vspace{0.5cm}
  \noindent{\color{lpAccent}\rule{\linewidth}{0.6pt}}
\end{titlepage}

\section*{Hinweise zu den Musterlösungen}
\thispagestyle{plain}

\noindent Eine mögliche Lösung. Mehrere Begründungswege sind legitim, solange sie:

\begin{itemize}
  \item den Begriffsapparat von Tag 6 nutzen (Omnichannel, Brüche, KPIs, Datenarchitektur, Operationalisierung),
  \item Annahmen sichtbar machen,
  \item Zielkonflikte (Tool vs.\ Organisation, Effizienz vs.\ Vollständigkeit) benennen,
  \item zu einer klaren Priorisierungsentscheidung führen.
\end{itemize}

\clearpage

\section*{Niveau L1 \textendash{} Wissen \& Reproduktion}
\addcontentsline{toc}{section}{L1 -- Wissen}

% --- AUFGABE 1 -----------------------------------------------------------
% LERNPLATT_HILFSVIDEOS
\section*{Hilfsvideos zu diesem Tag}
\addcontentsline{toc}{section}{Hilfsvideos zu diesem Tag}
\thispagestyle{plain}

\noindent Die folgenden YouTube-Erkl\"arvideos vermitteln die Inhalte, die Sie zur Bearbeitung der Aufgaben ben\"otigen. Klicken Sie auf den Titel, um das Video direkt zu \"offnen; die vollst\"andige URL steht in Klammern.

\begin{description}[style=nextline,leftmargin=18pt,labelindent=0pt,itemsep=8pt]
  \item[1.~Multichannel versus Omnichannel — der zentrale Unterschied] \href{https://youtu.be/pEGWLV6ZMOM}{\textcolor{lpAccent}{https://youtu.be/pEGWLV6ZMOM}}\\[2pt]
  {\footnotesize\color{lpGray}(URL: \nolinkurl{https://youtu.be/pEGWLV6ZMOM})}
  \item[2.~Lead Scoring — vom CRM-Eintrag zum priorisierten Lead] \href{https://youtu.be/tpVQBQWuLt8}{\textcolor{lpAccent}{https://youtu.be/tpVQBQWuLt8}}\\[2pt]
  {\footnotesize\color{lpGray}(URL: \nolinkurl{https://youtu.be/tpVQBQWuLt8})}
  \item[3.~KPIs im digitalen Marketing — Definition und Beispiele] \href{https://youtu.be/eXgtseYZtL8}{\textcolor{lpAccent}{https://youtu.be/eXgtseYZtL8}}\\[2pt]
  {\footnotesize\color{lpGray}(URL: \nolinkurl{https://youtu.be/eXgtseYZtL8})}
  \item[4.~Omnichannel als Architekturfrage — Plattformökonomie verstehen] \href{https://youtu.be/3B5Rn_lrGpQ}{\textcolor{lpAccent}{\nolinkurl{https://youtu.be/3B5Rn_lrGpQ}}}\\[2pt]
  {\footnotesize\color{lpGray}(URL: \nolinkurl{https://youtu.be/3B5Rn_lrGpQ})}
\end{description}
\vspace{0.6em}
\noindent{\footnotesize\color{lpGray}\textit{Tipp:} \"Offnen Sie das Video, lassen Sie es im Hintergrund laufen und arbeiten Sie parallel an der Aufgabe.}

\clearpage

\aufgabenkopf{1}{Single \textendash{} Multi \textendash{} Cross \textendash{} Omnichannel}{\nivLi}{6\,min}

\ytvideo{KPIs im digitalen Marketing — Definition und Beispiele}{https://youtu.be/eXgtseYZtL8}

\begin{loesungbox}
\textbf{1. Definitionen:}
\begin{itemize}
  \item \emph{Single Channel:} Kunde wird über genau \emph{einen} Kanal angesprochen \textendash{} z.\,B.\ ausschließlich Außendienst.
  \item \emph{Multichannel:} mehrere unabhängig betriebene Kanäle. Kunde sieht jeden Kanal isoliert; Kanäle wissen nichts voneinander.
  \item \emph{Cross-Channel:} Kanäle reichen \emph{Informationen} aneinander weiter (z.\,B.\ Webformular leitet an Vertrieb), aber Erlebnis ist noch nicht durchgängig konsistent.
  \item \emph{Omnichannel:} integriertes Kundenerlebnis. Konsistente Informationen, Preise, Ansprechpartner, Daten über alle Berührungspunkte hinweg.
\end{itemize}

\textbf{2. B2B-Beispiele:}
\begin{itemize}
  \item Single: traditioneller Maschinenbauer mit reinem Außendienstvertrieb.
  \item Multi: Mittelständler mit Website, LinkedIn und Außendienst, aber keine Verbindung zwischen ihnen.
  \item Cross: Website-Anfrage geht per E-Mail an Vertrieb, der die Anfrage manuell ins CRM kopiert.
  \item Omni: Lead über Webinar $\rightarrow$ automatisch in CRM $\rightarrow$ Lead-Score $\rightarrow$ Innendienst kontaktiert binnen 2\,Tagen $\rightarrow$ Außendienst sieht komplette Historie.
\end{itemize}

\textbf{3. Falsch verwendet:}
``Omnichannel'' wird in Marketingfolien am häufigsten falsch verwendet \textendash{} meist ist tatsächlich nur \emph{Multichannel} gemeint (``wir sind auf allen Kanälen präsent''). Echte Omnichannel-Integration erfordert technische, organisatorische und datenseitige Integration, die im Mittelstand selten erreicht wird.
\end{loesungbox}

\clearpage

% --- AUFGABE 2 -----------------------------------------------------------
\aufgabenkopf{2}{Fünf typische Omnichannel-Brüche}{\nivLi}{6\,min}

\ytvideo{Multichannel versus Omnichannel — der zentrale Unterschied}{https://youtu.be/pEGWLV6ZMOM}

\begin{loesungbox}
\textbf{1. Brüche + Beispiel:}
\begin{itemize}
  \item \emph{Doppelte Datenerfassung:} Kunde gibt Daten auf Website ein, der Außendienst fragt im Erstgespräch alles nochmal ab.
  \item \emph{Unklare Zuständigkeiten:} Webinar-Teilnehmende werden weder vom Marketing noch vom Vertrieb nachverfolgt.
  \item \emph{Widersprüchliche Informationen:} Website nennt 4 Wochen Lieferzeit, Vertrieb sagt 8 Wochen.
  \item \emph{Medienbrüche:} Anfrage per E-Mail $\rightarrow$ Vertrieb antwortet, übergibt aber nicht ins CRM.
  \item \emph{Fehlende Leadübergabe:} Marketing erfasst Lead, Vertrieb erfährt nichts oder erst Tage später.
\end{itemize}

\textbf{2. Schmerz \& Kosten:}
Für den \emph{Kunden} am ärgerlichsten: \emph{widersprüchliche Informationen}. Sie zerstören Vertrauen sofort. Für das \emph{Unternehmen} am teuersten: \emph{fehlende Leadübergabe}, weil daraus direkt verlorener Umsatz entsteht (kalter Lead nach 5 Tagen ist meist verloren).

\textbf{3. Sofort-Maßnahmen:}
\begin{itemize}
  \item Doppel-Erfassung: Vertrieb sieht Web-Daten vor dem Erstgespräch.
  \item Zuständigkeit: ein:e Eigentümer:in pro Kanal benennen.
  \item Widersprüche: zentrale ``Source of Truth''-Datenbank für Preise und Lieferzeit.
  \item Medienbrüche: CRM-Pflichtprozess für jede digitale Anfrage.
  \item Leadübergabe: Lead-SLA mit Reaktionszeit, automatischer Eskalation.
\end{itemize}
\end{loesungbox}

\clearpage

% --- AUFGABE 3 -----------------------------------------------------------
\aufgabenkopf{3}{Acht Operations-KPIs unterscheiden}{\nivLi}{8\,min}

\ytvideo{Multichannel versus Omnichannel — der zentrale Unterschied}{https://youtu.be/pEGWLV6ZMOM}

\begin{loesungbox}
\textbf{1. KPI-Zuordnung:}

\smallskip
\begin{tabularx}{\linewidth}{@{}lX@{}}
\toprule
\textbf{Steuerungsebene} & \textbf{KPI} \\
\midrule
Operativ                & Reaktionszeit, Conversion Rate, Kosten pro Lead, Umsatz pro Kanal \\
\midrule
Management              & Leadqualität, Abschlussquote, Wiederkaufrate, Kundenzufriedenheit, Kanalbeitrag \\
\midrule
Strategisch             & CLV, CAC vs.\ CLV-Verhältnis \\
\bottomrule
\end{tabularx}

\smallskip
\textbf{2. Gemeinsame KPIs Marketing + Vertrieb:}
\emph{Leadqualität} (Conversion MQL\,$\rightarrow$\,SQL) und \emph{Abschlussquote} (SQL\,$\rightarrow$\,Auftrag). Beide verbinden Marketing-Output und Vertriebs-Erfolg und können nicht von einer Seite allein optimiert werden.

\textbf{3. Im Mittelstand selten berichtet:}
\emph{Reaktionszeit auf digitale Anfragen}. Sie ist die wichtigste Frühwarn-KPI \textendash{} wenn sie kippt, kippen Vertrauen und Conversion innerhalb weniger Wochen. Wird selten gemessen, weil Verantwortung zwischen Marketing, Innendienst und Vertrieb diffus ist.

\begin{merkbox}[Managementhinweis]
\emph{Reaktionszeit ist die unterschätzteste KPI im B2B.} Sie wird selten berichtet, aber 5-Minuten-Reaktion vs.\ 2-Tage-Reaktion verändert die Abschlusswahrscheinlichkeit um Faktor 3\,--\,7.
\end{merkbox}
\end{loesungbox}

\clearpage

\section*{Niveau L2 \textendash{} Anwendung \& Transfer}
\addcontentsline{toc}{section}{L2 -- Anwendung}

% --- AUFGABE 4 -----------------------------------------------------------
\aufgabenkopf{4}{FitOffice Solutions \textendash{} Brüche erkennen}{\nivLii}{25\,min}

\ytvideo{Multichannel versus Omnichannel — der zentrale Unterschied}{https://youtu.be/pEGWLV6ZMOM}

\begin{loesungbox}
\textbf{1. \& 2. Brüche + Zuordnung + Auswirkung:}

\smallskip
\begin{tabularx}{\linewidth}{@{}lXX@{}}
\toprule
\textbf{Bruch} & \textbf{Bereich} & \textbf{Auswirkung} \\
\midrule
Website-Anfragen per E-Mail an Vertrieb         & Prozess + Daten   & Anfragen verloren / vergessen / nicht messbar \\
\midrule
Webinar-Teilnehmende nicht systematisch verfolgt & Verantwortung    & Hochwertige Leads gehen verloren \\
\midrule
Außendienst kennt Kampagnen oft nicht           & Kommunikation     & Lead-Gespräche ohne Kontext, schlechtere Conversion \\
\midrule
Marketing misst Klicks, Vertrieb Abschlüsse     & Verantwortung    & Konflikt, kein gemeinsames Ziel \\
\midrule
Kundendaten in Excel/CRM/E-Mail verteilt        & Daten            & Keine Einheitssicht; widersprüchliche Kommunikation \\
\midrule
Unterschiedliche Informationen je Kanal         & Kunde            & Vertrauensverlust, Verwirrung \\
\bottomrule
\end{tabularx}

\smallskip
\textbf{3. Drei Sofort-Maßnahmen:}
\begin{itemize}
  \item \emph{CRM-Pflichtprozess:} jede Website-/Webinar-/LinkedIn-Anfrage wird automatisch im CRM erfasst, mit Quelle, Status und Verantwortlichkeit.
  \item \emph{Webinar-Nachverfolgung:} 24\,h nach Webinar erhält jeder Teilnehmer eine personalisierte Folge-Mail; SAL-Übergabe an Vertrieb nach Lead-Score.
  \item \emph{Briefing Außendienst:} wöchentliches 30-Min-Update über laufende Kampagnen und ihre Botschaften.
\end{itemize}

\textbf{4. Welcher zuerst?}
\emph{CRM-Pflichtprozess.} Ohne ein zentrales System ist keine andere Maßnahme messbar oder skalierbar \textendash{} alle anderen Brüche sind \emph{Folgen} dieses einen. Wer Webinar-Nachverfolgung oder Lead Scoring ohne sauberes CRM einführt, baut auf Sand.
\end{loesungbox}

\clearpage

% --- AUFGABE 5 -----------------------------------------------------------
\aufgabenkopf{5}{FitOffice \textendash{} sieben Maßnahmen priorisieren}{\nivLii}{30\,min}

\ytvideo{Lead Scoring — vom CRM-Eintrag zum priorisierten Lead}{https://youtu.be/tpVQBQWuLt8}

\begin{loesungbox}
\textbf{1. Bewertung (1\,--\,5, 5\,=\,stärkster Wert):}

\smallskip
\begin{tabularx}{\linewidth}{@{}lccccc@{}}
\toprule
& \textbf{Wirk.} & \textbf{Aufw.} & \textbf{Kost.} & \textbf{Geschw.} & \textbf{Strat.} \\
\midrule
A: CRM bereinigen        & 5 & 3 & 3 & 3 & 5 \\
B: Lead Scoring          & 4 & 3 & 3 & 4 & 4 \\
C: Website-Anfragen      & 4 & 3 & 4 & 4 & 4 \\
D: Webinar-Nachverfolg.  & 4 & 2 & 2 & 5 & 3 \\
E: LinkedIn-Kampagnen    & 3 & 3 & 3 & 4 & 3 \\
F: Kundenportal          & 3 & 5 & 5 & 1 & 4 \\
G: Gemeinsame KPIs       & 5 & 1 & 1 & 5 & 5 \\
\bottomrule
\end{tabularx}

\smallskip
\textbf{2. Priorisierung:}
\begin{itemize}
  \item \emph{Sofort starten:} G (gemeinsame KPIs \textendash{} kostet nichts, wirkt sofort), A (CRM bereinigen), D (Webinar-Nachverfolgung).
  \item \emph{Später:} B (Lead Scoring nach A), C (Website-Anfragen optimieren), E (LinkedIn).
  \item \emph{Zunächst nicht:} F (Kundenportal \textendash{} zu komplex, falsches Zeitfenster).
\end{itemize}

\textbf{3. 12-Monats-Roadmap:}
\begin{itemize}
  \item \emph{Monat 1\,--\,2:} G + A (gemeinsame KPIs, CRM-Datenqualität).
  \item \emph{Monat 3\,--\,4:} D (Webinar-Nachverfolgung), B (Lead Scoring Pilot).
  \item \emph{Monat 5\,--\,7:} C (Website-Relaunch mit klaren Anfragepfaden).
  \item \emph{Monat 8\,--\,10:} E (LinkedIn-Kampagnen skalieren).
  \item \emph{Monat 11\,--\,12:} Pilot Kundenportal vorbereiten (nicht starten).
\end{itemize}

\textbf{4. Fünf KPIs:}
\begin{itemize}
  \item Reaktionszeit auf digitale Anfragen (Stunden).
  \item Conversion MQL\,$\rightarrow$\,SQL (\%).
  \item Conversion SQL\,$\rightarrow$\,Auftrag (\%).
  \item Anteil CRM-vollständiger Leads (\%).
  \item Kosten pro SQL (\euro).
\end{itemize}

\textbf{5. Zurückgestellt:}
\emph{Kundenportal (F)}. Trotz strategischer Attraktivität wird es im ersten Jahr nicht gebaut. Ohne CRM-Qualität, klare Verantwortlichkeiten und Datenarchitektur wird ein Portal zur teuren Insellösung mit unklarem Nutzen. Erst Fundament bauen.

\begin{merkbox}[Managementhinweis]
\emph{Maßnahme G (gemeinsame KPIs) kostet nichts und ist die wirksamste Einzelmaßnahme.} Trotzdem wird sie in 80\,\% der Mittelstandsprojekte zuletzt umgesetzt \textendash{} weil sie politisch unangenehm ist. Sie muss zuerst kommen, weil alle anderen Maßnahmen ohne sie ins Leere laufen.
\end{merkbox}
\end{loesungbox}

\clearpage

% --- AUFGABE 6 -----------------------------------------------------------
\aufgabenkopf{6}{Sechs Rollen in einer integrierten Vertriebsarchitektur}{\nivLii}{15\,min}

\ytvideo{KPIs im digitalen Marketing — Definition und Beispiele}{https://youtu.be/eXgtseYZtL8}

\begin{loesungbox}
\textbf{1. Hauptverantwortungen:}
\begin{itemize}
  \item \emph{Vertrieb:} SQL\,$\rightarrow$\,Auftrag, Bestandskundenpflege, Schlüsselkundengeschäft.
  \item \emph{Marketing:} MQL-Generierung, Content, Markenpositionierung, Kampagnen.
  \item \emph{IT:} CRM, Schnittstellen, Performance, Tracking.
  \item \emph{Kundenservice:} Erstkontakt, Reklamationen, Bestandskundenkommunikation.
  \item \emph{Datenmanagement:} Datenarchitektur, Reporting, Datenqualität.
  \item \emph{Management:} Strategie, KPI-Zielsetzung, Budget, Konfliktlösung.
\end{itemize}

\textbf{2. Reibungspunkte:}
\begin{itemize}
  \item Vertrieb $\leftrightarrow$ Marketing: Lead-Qualität.
  \item Marketing $\leftrightarrow$ IT: Tracking, Marketing-Tech-Stack.
  \item IT $\leftrightarrow$ Datenmanagement: wer entscheidet über Datenmodelle.
  \item Kundenservice $\leftrightarrow$ Vertrieb: Wo endet Service, wo beginnt Upsell?
  \item Management $\leftrightarrow$ Vertrieb: Quartalsziele vs.\ langfristige Strategie.
\end{itemize}

\textbf{3. Nicht eigenständig im Mittelstand:}
\emph{Datenmanagement}. Wird oft als Teilaufgabe der IT oder des Marketings behandelt. Folge: keine durchgängige Datenarchitektur, jedes Tool baut sein eigenes Modell, Reports widersprechen sich. Eigentumsfrage bleibt strukturell ungelöst.

\textbf{4. Steuerungskreis Digital Sales:}
\begin{itemize}
  \item \emph{Mitglieder:} CRO/Vertriebsleitung, Marketingleitung, IT-/Datenleitung, Kundenservice-Leitung.
  \item \emph{Frequenz:} monatlich 90\,Min., quartalsweise strategisch 4\,h.
  \item \emph{Entscheidungsrechte:} Budgetverteilung digitale Vertriebsmaßnahmen, KPI-Definitionen, Lead-Übergabe-Regeln, Konfliktlösung an Marketing/Vertrieb-Schnittstelle.
\end{itemize}
\end{loesungbox}

\clearpage

% --- AUFGABE 7 -----------------------------------------------------------
\aufgabenkopf{7}{IndustrialPro Services \textendash{} Omnichannel-Architektur}{\nivLii}{30\,min}

\ytvideo{Lead Scoring — vom CRM-Eintrag zum priorisierten Lead}{https://youtu.be/tpVQBQWuLt8}

\begin{loesungbox}
\begin{center}
\begin{tikzpicture}[font=\sffamily\footnotesize,
  hub/.style={circle, draw=lpAccent, fill=lpAccent, text=white,
    font=\sffamily\bfseries\small, minimum size=1.8cm, align=center},
  ch/.style={rectangle, rounded corners=3pt, draw=lpGray!70, fill=lpGrayLight,
    text=lpInk, minimum width=1.9cm, minimum height=0.8cm, align=center,
    font=\sffamily\footnotesize},
  prio/.style={ch, draw=lpAmberBorder, fill=lpAmberLight, line width=0.8pt},
  flow/.style={-{Stealth[length=4pt]}, line width=0.7pt, color=lpAccent!70}]
  % CRM-Hub im Zentrum
  \node[hub] (crm) at (0,0) {CRM\\(Hub)};
  % Kanäle ringsum
  \node[prio]  (web)  at ( 3.3, 1.6) {Website};
  \node[prio]  (lead) at ( 3.3,-1.6) {Lead\\Scoring};
  \node[ch]    (mail) at ( 0.0, 2.5) {E-Mail};
  \node[ch]    (port) at ( 0.0,-2.5) {Kunden-\\portal};
  \node[ch]    (ad)   at (-3.3, 1.6) {Außen-\\dienst};
  \node[prio]  (dash) at (-3.3,-1.6) {Dashboard};
  % Verbindungen (bidirektional)
  \foreach \n in {web,lead,mail,port,ad,dash} {
    \draw[flow] (crm) -- (\n);
    \draw[flow] (\n) -- (crm);
  }
  % Legende
  \node[font=\sffamily\scriptsize\itshape, color=lpAmberBorder] at (0,-3.4)
    {Orange = Priorit{\"a}t 1--4 (CRM, Lead-Scoring, Website, Dashboard) \textbullet{} Grau = sp{\"a}ter / Pilot};
\end{tikzpicture}
\end{center}

\smallskip
\textbf{1. Omnichannel-Schwächen:}
\begin{itemize}
  \item Bestehende Kanäle isoliert \textendash{} keine integrierte Strategie.
  \item Hauptproblem ist nicht Kanalpräsenz, sondern Prozess- und Datenintegration.
  \item CRM nicht als zentrale Steuerung genutzt.
  \item Marketing, Vertrieb, Kundenservice mit verschiedenen Zielgrößen.
  \item Kunden erleben Medienbrüche und uneinheitliche Kommunikation.
\end{itemize}

\textbf{2. Ideale Customer Journey für Service-Lead:}
Kunde erkennt Wartungs-/Modernisierungsbedarf $\rightarrow$ recherchiert online (Referenzen, Reaktionszeiten) $\rightarrow$ stellt Serviceanfrage über Website oder Kundenportal $\rightarrow$ Anfrage automatisch ins CRM erfasst und qualifiziert $\rightarrow$ Lead Scoring entscheidet: Vertrieb / Kundenservice / Außendienst $\rightarrow$ Rückmeldung innerhalb definierter Zeit $\rightarrow$ Angebot/Termin dokumentiert $\rightarrow$ nach Abschluss in Wartungs-/Folgeprozess überführt.

\textbf{3. Maßnahmenbewertung (zusammengefasst):}
\begin{itemize}
  \item \emph{A CRM-Pflichtprozess:} höchste Priorität \textendash{} ohne CRM keine Steuerung.
  \item \emph{B Lead Scoring + Übergaberegeln:} sehr wichtig.
  \item \emph{C Website-Relaunch:} wichtig, fokussiert auf klare Leadpfade.
  \item \emph{D E-Mail-Strecken:} sinnvoll nach sauberer Datenbasis.
  \item \emph{E Kundenportal:} strategisch wertvoll, aber komplex \textendash{} als Pilot starten.
  \item \emph{F Dashboard:} wichtig für gemeinsame Führungssicht.
  \item \emph{G Außendienst-Schulung:} notwendig für Akzeptanz.
\end{itemize}

\textbf{4. Priorisierung:}
\begin{itemize}
  \item Priorität 1: A (CRM-Pflichtprozess + Datenqualität).
  \item Priorität 2: B (Lead Scoring + Übergaberegeln).
  \item Priorität 3: C (Website mit digitalen Anfragewegen).
  \item Priorität 4: F (gemeinsames Dashboard).
  \item Priorität 5: G (Außendienst-Schulung).
  \item E (Kundenportal) als \emph{Pilot} mit Bestandskunden, nicht als Großprojekt.
\end{itemize}

\textbf{5. Roadmap mit Verantwortlichkeiten:}
\begin{itemize}
  \item \emph{Monat 1\,--\,3:} A + Lead-Status-Definition (CRO + IT).
  \item \emph{Monat 4\,--\,6:} B + Übergaberegeln (CRO + Marketingleitung).
  \item \emph{Monat 7\,--\,9:} C + F (Marketing + IT).
  \item \emph{Monat 10\,--\,12:} G + E-Pilot (Vertriebsleitung + Service).
\end{itemize}

\begin{merkbox}[Managementhinweis]
\emph{Omnichannel beginnt bei Daten und Prozessen, nicht bei Tools.} Wer mit einem teuren Marketing-Tool startet, ohne CRM-Pflichtprozess und Übergaberegeln zu klären, hat nach 12\,Monaten teure Tools und dieselben Brüche.
\end{merkbox}
\end{loesungbox}

\clearpage

% --- AUFGABE 8 -----------------------------------------------------------
\aufgabenkopf{8}{Datenarchitektur \& Single Source of Truth}{\nivLii}{20\,min}

\ytvideo{KPIs im digitalen Marketing — Definition und Beispiele}{https://youtu.be/eXgtseYZtL8}

\begin{loesungbox}
\textbf{1. Warum verteilte Daten strategisch problematisch sind:}
\begin{itemize}
  \item \emph{Widersprüche:} jede Quelle ``hat recht'' \textendash{} Reports erzeugen Misstrauen.
  \item \emph{Skalierbarkeit:} ohne SSoT skaliert keine Personalisierung, kein Lead Scoring, keine Wiederkauflogik.
  \item \emph{Lock-in:} Daten in verteilten Systemen sind faktisch \emph{nicht mehr eigen} \textendash{} sie liegen bei Tool-Anbietern.
  \item \emph{Vertrieb-Marketing-Konflikt:} ohne SSoT bleibt der Lead-Streit politisch.
\end{itemize}

\textbf{2. Fünf Pflicht-CRM-Felder:}
\begin{itemize}
  \item Kunden-ID + Status (Lead, MQL, SAL, SQL, Bestand).
  \item Quelle (Webinar, LinkedIn, Empfehlung, etc.).
  \item Letzter Kontakt + Verantwortliche:r.
  \item Bedarfstyp und Zeitfenster (geschätzt).
  \item Kommunikationshistorie (alle Berührungspunkte).
\end{itemize}

\textbf{3. Wer entscheidet \textendash{} und warum oft falsch:}
Im Mittelstand entscheidet typischerweise die \emph{IT-Leitung} über Datenarchitektur. Das ist oft die falsche Person, weil die Architektur-Frage primär \emph{geschäftlich} ist (welche Daten brauchen wir, warum?), nicht technisch. Die Entscheidung sollte beim CRO oder einer Daten-/Customer-Insights-Funktion liegen, die IT als Umsetzer einbindet.

\textbf{4. Datenfluss digitaler Lead:}
Website-Formular ausgefüllt $\rightarrow$ \emph{automatisch} ins CRM (kein E-Mail-Detour) $\rightarrow$ Lead-Score berechnet (basiert auf Pflicht-Feldern + Aktivitätsdaten) $\rightarrow$ bei Score $\geq$\,X: Push an Innendienst (24\,h-SLA) $\rightarrow$ Innendienst qualifiziert, ändert Status auf SAL/SQL $\rightarrow$ bei SQL: Übergabe an Außendienst (5\,Werktage-SLA) $\rightarrow$ jeder Schritt im CRM dokumentiert.

\begin{merkbox}[Lessons Learned]
\emph{Datenarchitektur ist eine Geschäftsentscheidung, keine IT-Entscheidung.} Wer sie der IT überlässt, bekommt eine technisch saubere, aber geschäftlich wertlose Architektur.
\end{merkbox}
\end{loesungbox}

\clearpage

\section*{Niveau L3 \textendash{} Management \& Entscheidungsarchitektur}
\addcontentsline{toc}{section}{L3 -- Management}

% --- AUFGABE 9 -----------------------------------------------------------
\aufgabenkopf{9}{IndustrialPro \textendash{} 12-Monats-Omnichannel-Entscheidung}{\nivLiii}{45\,min}

\ytvideo{Multichannel versus Omnichannel — der zentrale Unterschied}{https://youtu.be/pEGWLV6ZMOM}

\begin{loesungbox}
\textbf{1. Priorisierte Maßnahmenentscheidung (5 Hebel):}
\textbf{A} CRM-Pflichtprozess, \textbf{B} Lead Scoring + Übergaberegeln, \textbf{C} Website-Relaunch, \textbf{F} gemeinsames Dashboard, \textbf{G} Außendienst-Schulung. Begründung: Diese fünf Maßnahmen bilden zusammen ein \emph{Vertriebsbetriebssystem}. CRM (A) ist die Datenbasis; B definiert die Übergabe; C macht die Website zum Conversion-Hub; F gibt die gemeinsame Führungssicht; G sichert die Vertriebs-Akzeptanz. Maßnahme D (E-Mail) wird zurückgestellt, weil sie ohne saubere Datenbasis Lead-Müll skaliert. Maßnahme E (Kundenportal) wird zum Pilot mit 30 Bestandskunden, nicht zum Großprojekt.

\textbf{2. Budget (250.000\,\euro):}
\begin{itemize}
  \item 70.000\,\euro{} \textendash{} CRM-Konsolidierung (A): Datenbereinigung, Pflichtfelder, automatische Lead-Erfassung.
  \item 40.000\,\euro{} \textendash{} Lead Scoring + Übergaberegeln (B).
  \item 60.000\,\euro{} \textendash{} Website-Relaunch (C): Serviceanfrage, Beratungstermin, Referenzen.
  \item 30.000\,\euro{} \textendash{} Dashboard (F).
  \item 20.000\,\euro{} \textendash{} Außendienst-Schulung (G).
  \item 15.000\,\euro{} \textendash{} Pilot Kundenportal mit 30 Wartungsvertrags-Kunden.
  \item 15.000\,\euro{} \textendash{} Tests, KPI-Setup, externe Verstärkung.
\end{itemize}

\textbf{3. 12-Monats-Roadmap in 4 Phasen:}

\emph{Monat 1\,--\,3:} CRM-Pflichtfelder definiert, gemeinsame KPIs zwischen Marketing/Vertrieb/Service vereinbart, Lead-SLA festgelegt.

\emph{Monat 4\,--\,6:} Lead Scoring eingeführt, Website-Anfrageprozesse verbessert, Außendienst-Schulungen durchgeführt.

\emph{Monat 7\,--\,9:} Dashboard live, E-Mail-Nurturing pilotiert (intern, ohne Außenwirkung), erste Auswertung Lead-Qualität.

\emph{Monat 10\,--\,12:} Kundenportal-Pilot mit 30 Bestandskunden, Prozessoptimierung, Skalierungsentscheidung für Kundenportal.

\textbf{4. Sieben KPIs entlang Funnel:}
\begin{itemize}
  \item Anzahl digitaler Leads (Reichweite).
  \item Reaktionszeit auf digitale Anfrage (h).
  \item Conversion MQL\,$\rightarrow$\,SQL (\%).
  \item Conversion SQL\,$\rightarrow$\,Angebot (\%).
  \item Abschlussquote digitaler Leads (\%).
  \item Anteil CRM-vollständiger Leads (\%).
  \item Kosten pro qualifiziertem Lead (\euro).
\end{itemize}

\textbf{5. Governance:}
\begin{itemize}
  \item CRO entscheidet über digitale Vertriebsstrategie, Budget, Konflikte.
  \item Monatlicher Steuerungskreis ``Digital Sales'' (CRO, Marketingleitung, Vertriebsleitung, IT, Kundenservice).
  \item Quartalsweiser Strategie-Review (Vorstand).
  \item Wöchentliche operative Stand-ups (Marketing + Vertrieb).
\end{itemize}

\textbf{6. Entscheidung gegen attraktive Maßnahme:}
\emph{Kundenportal nicht als Großprojekt im Jahr 1.} Trotz strategischer Bedeutung und Investorenpush wird das Kundenportal nur als Pilot mit 30 Wartungsvertrags-Kunden umgesetzt. Begründung: ohne stabile CRM-Prozesse, klare Verantwortlichkeiten und Datenqualität würde das Portal zusätzliche Komplexität erzeugen \textendash{} und scheitern. Der Pilot testet Kundennutzen und Prozessfähigkeit kontrolliert.

\textbf{7. Entscheidung unter Unsicherheit:}
Investition von 60.000\,\euro{} in Website-Relaunch \emph{vor} klarem Beweis, dass digitale Anfragen die Hauptquelle für Servicegeschäft werden. Annahme: \emph{Die Customer Journey für Service-Leads beginnt online; Anbieter ohne digitale Anfragepfade fallen aus dem Vergleich.} Lernindikator: digital generierte Service-Anfragen $\geq$\,80 pro Monat bis Monat 9. Wenn nicht: Re-Allokation in andere Hebel.

\begin{merkbox}[Senior-Level-Hinweis]
\emph{Operationalisierung ist Sequenz-Disziplin, nicht Maßnahmenbreite.} Wer fünf Hebel in 12\,Monaten konsequent umsetzt, gewinnt; wer zehn parallel startet, verliert. Senior-Level-Disziplin heißt: \emph{aktiv etwas nicht zu machen.}
\end{merkbox}
\end{loesungbox}

\clearpage

% --- AUFGABE 10 ----------------------------------------------------------
\aufgabenkopf{10}{Operationalisierung als Organisationsentwicklung}{\nivLiii}{30\,min}

\ytvideo{Multichannel versus Omnichannel — der zentrale Unterschied}{https://youtu.be/pEGWLV6ZMOM}

\begin{loesungbox}
\textbf{1. Warum Omnichannel ohne Organisationsentwicklung scheitert:}
Omnichannel ist primär ein \emph{Schnittstellen-} und \emph{Anreizproblem}. Tools können Daten verbinden, aber nicht Menschen mit unterschiedlichen Zielen. Solange Marketing auf Reichweite, Vertrieb auf Abschlüsse, Kundenservice auf Tickets gemessen wird, bauen die drei systematisch \emph{gegen} Omnichannel-Logik. Ohne Anreiz-, Rollen- und KPI-Anpassung produzieren Tools nur teurer dokumentierte Brüche.

\textbf{2. Drei Anti-Pattern:}
\begin{itemize}
  \item \emph{``Wir kaufen ein neues CRM.''} Neues System, alte Datenchaos, alte Anreize. Effekt: 6\,Monate Migration, alle Probleme zurück.
  \item \emph{``Marketing macht jetzt Lead Gen, Vertrieb macht Abschlüsse.''} Saubere Trennung klingt logisch, ignoriert aber die Übergabe \textendash{} dort liegt der Engpass.
  \item \emph{``Wir schulen den Außendienst auf digitale Tools.''} Schulung ohne Anreizanpassung verändert kein Verhalten. Außendienst macht weiter, was provisioniert wird.
\end{itemize}

\textbf{3. Organisationsentwicklungs-Konzept IndustrialPro:}
\begin{itemize}
  \item \emph{Rollen:} CRO als digitales Vertriebsleitungs-Mandat; Innendienst als Lead-Qualifizierungsfunktion; Marketing mit Lead-Qualitäts-Verantwortung.
  \item \emph{Anreize:} Provisionsmodell Außendienst um Anteil ``digitaler Lead-Conversion'' ergänzen. Marketing-Boni teilweise an SQL-Quote koppeln.
  \item \emph{Gemeinsame KPIs:} MQL\,$\rightarrow$\,SQL-Conversion (Marketing + Vertrieb), SQL\,$\rightarrow$\,Auftrag (Vertrieb).
  \item \emph{Steuerungskreise:} wöchentlich operativ, monatlich CRO-Level, quartalsweise Vorstand.
  \item \emph{Kompetenzaufbau:} digitale Verkaufstechnik für Außendienst, Lead-Coaching für Innendienst, Data Literacy für Marketing.
\end{itemize}

\textbf{4. Wirksamste Einzelmaßnahme:}
\emph{Provisionsmodell-Anpassung beim Außendienst.} Solange der Außendienst nicht von digitalen Leads profitiert, sabotiert er sie strukturell. Eine Provisionsanpassung verändert Verhalten in Wochen \textendash{} während Schulungen Monate brauchen und nur teilweise wirken.

\textbf{5. Faustregel \textendash{} startbereit:}
Ein Omnichannel-Projekt ist organisatorisch startbereit, wenn (a)~gemeinsame KPIs definiert sind, (b)~das Provisionsmodell die Anreize ausgerichtet hat, und (c)~ein:e Eigentümer:in der Customer Journey benannt ist. Ohne diese drei ist jeder weitere Schritt vorzeitig.

\begin{merkbox}[Lessons Learned]
\emph{Tools verändern keine Organisation. Anreize tun es.} Wer Omnichannel ohne Anreizdiskussion startet, hat in 12\,Monaten teuere Tools und alte Konflikte.
\end{merkbox}
\end{loesungbox}

\clearpage

% --- AUFGABE 11 ----------------------------------------------------------
\aufgabenkopf{11}{Transferprojekt \textendash{} Omnichannel-Vertriebsarchitektur}{\nivLiii}{45\,min}

\ytvideo{Multichannel versus Omnichannel — der zentrale Unterschied}{https://youtu.be/pEGWLV6ZMOM}

\begin{loesungbox}
\emph{Musterlösung als Entscheidungsarchitektur \textendash{} andere Konfigurationen möglich.}

\smallskip
\textbf{1. Zielbild:}
Das Unternehmen betreibt in 18\,Monaten eine integrierte Omnichannel-Architektur, in der jede Kundenanfrage \textendash{} digital oder persönlich \textendash{} in einem System der Wahrheit landet, klar zugeordnet wird und mit definierten Reaktionszeiten bearbeitet wird. Marketing und Vertrieb arbeiten mit gemeinsamen KPIs; Kunden erleben konsistente Information über alle Kanäle. Kunden-Daten gehören dem Unternehmen, nicht den Tool-Anbietern. Reaktionszeit auf digitale Anfragen $<$\,4\,h, Conversion MQL\,$\rightarrow$\,SQL $>$\,25\,\%.

\textbf{2. Customer Journey:}
Aufmerksamkeit (LinkedIn, SEO, Empfehlung) $\rightarrow$ Informationssuche (Website, Whitepaper, Webinar) $\rightarrow$ Vergleich (Beratungstermin / Konfigurator) $\rightarrow$ Entscheidung (Außendienst + Angebot) $\rightarrow$ Kauf (CRM-dokumentiert) $\rightarrow$ Nutzung (Kundenportal / Service) $\rightarrow$ Wiederkauf (E-Mail-Automation + Innendienst).

\textbf{3. Kanalbrüche und Prozessrisiken (Analyse):}
\begin{itemize}
  \item Daten verteilt $\rightarrow$ keine SSoT.
  \item Lead-Übergabe undefiniert $\rightarrow$ kalte Leads.
  \item Marketing/Vertrieb mit verschiedenen KPIs.
  \item Außendienst nicht in digitale Kampagnen eingebunden.
  \item Kundenservice ohne CRM-Anbindung.
\end{itemize}

\textbf{4. Priorisierung Kanäle und Übergaben:}
\begin{itemize}
  \item Website mit Conversion-Logik (Hub).
  \item LinkedIn-Selling (Reichweite).
  \item Webinare (Leadqualifikation).
  \item Außendienst (Schlüsselkunden).
  \item Innendienst als Lead-Qualifizierer (zentrale Übergabestelle).
  \item Kundenportal (Bestandskunden, Pilot).
\end{itemize}

\textbf{5. Prozessarchitektur:}
\begin{itemize}
  \item Lead-Aufnahme: jede Anfrage automatisch ins CRM.
  \item Qualifizierung: Lead-Scoring nach Pflichtfeldern + Aktivität.
  \item Übergabe: Innendienst (MQL\,$\rightarrow$\,SAL, 24\,h), Außendienst (SAL\,$\rightarrow$\,SQL, 5\,Werktage).
  \item Nachverfolgung: alle Touchpoints im CRM dokumentiert.
  \item Abschluss: Auftragsdaten + Wiederkauf-Trigger gesetzt.
\end{itemize}

\textbf{6. Datenarchitektur:}
\begin{itemize}
  \item CRM als SSoT.
  \item Pflichtfelder: ID, Quelle, Status, Verantwortliche:r, Bedarfstyp, Zeitfenster, Kommunikationshistorie.
  \item Lead-Score, MQL/SAL/SQL/Bestandsstatus.
  \item Reporting: gemeinsames Dashboard für Marketing, Vertrieb, Kundenservice.
\end{itemize}

\textbf{7. Governance:}
\begin{itemize}
  \item CRO: Strategie + Budget.
  \item Marketingleitung: Lead-Qualität, Content.
  \item Vertriebsleitung: SQL\,$\rightarrow$\,Auftrag.
  \item IT/Daten: Architektur, Tracking, Reporting.
  \item Kundenservice: Erstkontakt, Reklamationen, Bestandskommunikation.
  \item Monatlicher Steuerungskreis als Entscheidungsorgan.
\end{itemize}

\textbf{8. KPI-Architektur:}
\begin{itemize}
  \item \emph{Operativ:} Reaktionszeit, Lead-Anzahl pro Quelle, CRM-Vollständigkeit.
  \item \emph{Management:} MQL\,$\rightarrow$\,SQL-Conversion, SQL\,$\rightarrow$\,Auftrag, Kosten pro SQL.
  \item \emph{Strategisch:} CLV, CAC vs.\ CLV, Anteil digitaler Neukunden.
\end{itemize}

\textbf{9. Roadmap 12\,--\,18 Monate:}
\begin{itemize}
  \item Monat 1\,--\,3: Gemeinsame KPIs, CRM-Pflichtfelder, Lead-SLA, Provisionsmodell anpassen.
  \item Monat 4\,--\,6: Lead Scoring + Website-Relaunch + Außendienst-Schulung.
  \item Monat 7\,--\,9: Dashboard, E-Mail-Automation, erste Auswertung.
  \item Monat 10\,--\,12: Kundenportal-Pilot mit 30 Kunden.
  \item Monat 13\,--\,18: Skalierung erfolgreicher Hebel, Datenarchitektur erweitern, Bestandskunden-Wiederkauf-Logik.
\end{itemize}

\textbf{10. Drei Managemententscheidungen unter Unsicherheit:}
\begin{enumerate}
  \item \emph{Provisionsmodell anpassen vor klarem Pilotnachweis.} Annahme: ohne Anreizanpassung kein Verhaltenswandel.
  \item \emph{Investition in CRM-Datenqualität vor sichtbarem ROI.} Annahme: ohne SSoT skaliert keine andere Maßnahme.
  \item \emph{Verzicht auf Kundenportal-Großprojekt im Jahr 1.} Annahme: ohne Fundament wird das Portal teuer und wirkungslos.
\end{enumerate}

\textbf{Zusatz \textendash{} Verzicht auf attraktive Initiative:}
Bewusster Verzicht auf eine zusätzliche Marketing-Automation-Suite, die Marketing im Jahr 1 eingeführt sehen möchte. Aus Senior-Level-Sicht: ohne saubere Datenbasis und gemeinsame KPIs erzeugt mehr Automation nur schneller Lead-Müll. Die ``verlorene'' Reichweite ist eine \emph{Investition in spätere Skalierbarkeit}.

\begin{merkbox}[Senior-Level-Hinweis]
\emph{Eine Omnichannel-Architektur ist Organisationsdesign \textendash{} nicht IT-Auswahl.} Wer Tools wählt, bevor Anreize und Rollen geklärt sind, kauft Software, nicht Wirkung.
\end{merkbox}
\end{loesungbox}

\clearpage

\section*{Abschluss}
\thispagestyle{plain}

\noindent Die Musterlösungen sind eine Diskussionsbasis.

\medskip
\begin{infobox}[Nächster Schritt]
Im Coaching: Welche Anti-Pattern (``wir kaufen ein neues CRM'', ``wir schulen den Außendienst'') haben Sie in Ihrer Organisation schon erlebt \textendash{} und welche Anreizanpassung würde sie strukturell verhindern?
\end{infobox}

\vspace{1cm}
\noindent{\color{lpAccent}\rule{\linewidth}{0.6pt}}\\[4pt]
{\footnotesize\sffamily\color{lpGray}Lernplatt \textperiodcentered{} by Can Siebert \textperiodcentered{} Kurs 22 (Bo 143) \textperiodcentered{} Tag 6 \textperiodcentered{} Lösungsheft \textperiodcentered{} \textcopyright{} 2026}

\end{document}
